\documentclass{article}
\usepackage{amsmath, amssymb, amsthm}
\usepackage{hyperref}
\usepackage{geometry}
\geometry{margin=1in}

\title{Erd\H{o}s Problem \#272}
\date{Last updated: 2025-08-31}
\author{Source: erdosproblems.com}

\begin{document}
\maketitle

\noindent\textbf{Status:} OPEN \quad \textbf{No prize}

\noindent\textbf{Tags:} additive combinatorics, arithmetic progressions

\medskip

\noindent\textbf{Problem Statement:}

Let $N\geq 1$. What is the largest $t$ such that there are $A_1,\ldots,A_t\subseteq \{1,\ldots,N\}$ with $A_i\cap A_j$ a non-empty arithmetic progression for all $i\neq j$?




Simonovits and S\'{o}s \cite{SiSo81} have shown that $t\ll N^2$. Erd\H{o}s and Graham asked whether the maximal $t$ is achieved when we take the $A_i$ to be all arithmetic progressions in $\{1,\ldots,N\}$ containing some fixed element, 'presumably the integer $\lfloor N/2\rfloor$'. This was disproved by Simonovits and S\'{o}s \cite{SiSo81}, who observed that taking all sets containing at most $3$ elements, containing some fixed element, produces $\binom{N}{2}+1$ many such sets, which is asymptotically greater than the number of arithmetic progressions containing a fixed element, which is $\sim \frac{\pi^2}{24}N^2$. If we drop the non-empty requirement then Graham, Simonovits, and S\'{o}s \cite{GSS80} have shown that\[t\leq \binom{N}{3}+\binom{N}{2}+\binom{N}{1}+1\]and this is best possible. Szabo \cite{Sz99} proved that the maximal such $t$ is equal to\[\frac{N^2}{2}+O(N^{5/3}(\log N)^3),\]resolving the asymptotic question. On the other hand, Szabo showed that the conjecture of Simonovits and S\'{o}s that $\binom{n}{2}+1$ is best possible is false, giving a construction which yields\[t \geq \binom{N}{2}+\left\lfloor\frac{N-1}{4}\right\rfloor+1.\]Szabo conjectures that the asymptotic $t=\binom{N}{2}+O(N)$ holds, and that in any extremal example there is an integer contained in all sets.




References

[GSS80] Graham, R. L. and Simonovits, M. and S\'{o}s, V. T., A note on the intersection properties of subsets of integers . J. Combin. Theory Ser. A (1980), 106-110.

[SiSo81] Simonovits, Mikl\'{o}s and S\'{o}s, Vera T., Intersection properties of subsets of integers . European J. Combin. (1981), 363-372.

[Sz99] Szab\'o, Tibor, Intersection properties of subsets of integers . European J. Combin. (1999), 429--444.

\noindent\textbf{Related OEIS sequences:} possible


\bigskip
\noindent\small{Source: \url{https://www.erdosproblems.com/272}}
\end{document}
